\begin{abstract}
Adaptation burden under drift is not determined by a single scalar notion of
memory quality. The empirical object in this paper is the minimum realizable
re-learning burden
\[
  R^*(m,\pi) \coloneqq \min_w R(m,\pi,w),
\]
defined over protocol-window choices after drift. The main empirical result is a
hierarchical decomposition. Across architectural families, effective spectral
rank separates degradation regimes. Within a recurrent local regime, activation
velocity $\Delta_v$ predicts burden strongly even when effective rank still
varies substantially. The pooled data therefore support a two-level structure:
rank is a family-level regime coordinate, while $\Delta_v$ is a within-regime
scale coordinate. This decomposition makes clear why static endpoint metrics are
insufficient: models with similar task performance can occupy different burden
regimes and require different post-drift re-learning effort.
\end{abstract}
